\documentclass[10pt]{amsart}
\usepackage[margin=1.6cm,top=1.5cm,bottom=1.7cm]{geometry}
\usepackage{amsmath,amssymb}
\let\olditemize\itemize\renewcommand{\itemize}{\olditemize\setlength{\itemsep}{1pt}\setlength{\topsep}{2pt}\setlength{\parsep}{0pt}}
\setlength{\parskip}{2pt}
\newcommand{\Z}{\mathbb Z}\newcommand{\Q}{\mathbb Q}\newcommand{\F}{\mathbb F}
\newcommand{\CP}{\mathbb{CP}}\newcommand{\PD}{\operatorname{PD}}
\newcommand{\Hol}{\operatorname{Hol}}
\title{Purely cosmetic surgeries via a $\mathrm{PU}(2)$-monopole cobordism:\\ a digest of the argument}
\date{}
\makeatletter\def\section{\@startsection{section}{1}{\z@}{.5\baselineskip}{.2\baselineskip}{\normalfont\bfseries}}\makeatother
\begin{document}
\maketitle
\vspace{-2.5em}

\noindent\textbf{Setting.} $Y_r=S^3_r(K)$. By Hanselman and
Daemi--Lidman--Miller Eismeier (Cor.~1.4), the only open case is $g(K)=2$
with an orientation-preserving diffeomorphism $\phi\colon Y_{-2}\to Y_2$.
Everything below derives a contradiction from $\phi$; no peripheral
condition on $\phi$ is used.

\noindent\textbf{Strategy.} Use $\phi$ to build a closed $4$-manifold $X$
with an $n$-cube of metrics and a signed sum over $2^n$ determinant lifts
such that (a)~the cut-down zero-dimensional instanton count is an integer
$\Omega\ne0$, and (b)~the coupled Dirac index $n_D=\Theta-\kappa$ is
positive. Run the $\mathrm{PU}(2)$-monopole cobordism
(Pidstrigach--Tyurin, Okonek--Teleman, Feehan--Leness): after the
$S^1$-quotient and $\eta=n_D-1$ phase cuts (sections of the weight-two
line) the free monopole space is $1$-dimensional. Its ends are instanton
links ($2^\eta\Omega$), Seiberg--Witten-type reducibles, and face ends.
Lattice and curvature estimates remove \emph{all} reducible fixed points
and all face ends except lens walls. The lens walls cancel in pairs with the
same signs that make $\Omega$ integral. Hence $2^\eta\Omega=0$. The analysis
never glues at a Seiberg--Witten point or at a bubbling level.

\section*{1. Ordinary maps (paper, Secs.~2--3)}
Use the irreducible determinant-one complexes. $Y_{\pm1},Y_{\pm2},S^3$ have
only central reducibles. On $Y_0$, $w\cdot A$ is odd ($A$ the capped Seifert
surface). The unperturbed complex gives $C(S^3)=0$.
\begin{itemize}
\item $f_+\colon Y_0\to Y_1\to Y_2$ (meridional $2$-handles) and its mirror
$f_-\colon Y_{-2}\to Y_0$ are $\F_2$-isomorphisms by exact-triangle
pentagons (Culler--Daemi--Xie, DLME); the third term vanishes because
$C(S^3)=0$. The outer cap is a sphere chain bounded by $S^1\times S^2$.
Replacing it by $S^1\times D^3$ gives continuation, and the harmonic path of
boundary flats has odd degree.
$g_+\colon Y_2\to S^3\to Y_0$ and $g_-$ are one-parameter interval counts,
summed mod~$2$ over two lifts differing by $\PD$ of the $(-2)$-sphere; the
$\mathbb{RP}^3$-wall terms agree and cancel mod~$2$. The compressed pentagon
gives $f_+g_+\simeq\mathrm{id}$.
\item \emph{Negative interval $B$ (first new ingredient).}
$W'=(-C_2)\cup_{S^3}C_{-2}\colon Y_2\to Y_{-2}$ has $b^+=0$ and
$H_2=\langle F_l,F_r\rangle$ with $F^2=-2$. The sphere is $S=F_l-F_r$, $S^2=-4$.
The metric interval runs from the $S^3$-split $J$ to splitting off
$N=\nu S$, $\partial N=L(4,1)$. The degree-two insertion $x(S)$ requires the
normalized holonomy $=-1$ on a one-parameter sweep of loops across $S$. The
count has $i+1+1-3=0$ and sums over lifts $c_e=e\PD(S)$. At the lens wall
only the minimal cap survives: trace-zero flat, $\kappa=\frac14$, harmonic
reducible. Its eigenline has degree one, so the sweep crosses $-1$ with
signed count $\pm1$. A local determinant-line comparison (tensor by the line
$-2x$ on $N$, then a boundary Weyl conjugation) gives weights
$\varepsilon_1=-s\varepsilon_0$ that cancel these terms \emph{integrally}.
Hence $B$ is an integral chain map. Mod~$2$, $B\simeq g_-g_+$: in the
$3$-dimensional associahedron for $2\to\infty\to0\to\infty\to-2$, the face
$\nu U$ ($U^2=0$) becomes $B$ in two steps. First replace $\nu U$ by
$S^1\times D^3$; the $w_2$ root mismatch makes the matching condition
$\Hol_\beta=-1$. Then contract $\beta$ through two discs whose union is
$\pm S$.
\item Bridge $T=f_-\phi^{-1}f_+\colon I(Y_0)\to I(Y_0)$ is a mod-$2$ unit,
hence a $\Q$-isomorphism.
\end{itemize}

\section*{2. Caps with a nonzero pairing on $Y_2$ (paper, Sec.~4)}
Take a symplectic closure $X_0\supset Y_0$ (Gabai, Eliashberg--Thurston,
Eliashberg caps, as in Kronheimer--Mrowka). Arrange simply connected halves,
tight genus-$2$ surfaces in each half, and $A$ primitive.
$D^w_{X_0}(x_0^mh^a)\ne0$ with $h\cdot A=1$; the paper reproves this. The
extremal joint eigenspace of $(\mu(\Sigma),\mu(\mathrm{pt}))$ on
$I(\Sigma_g\times S^1)$ is a line (Mu\~noz). The Lefschetz word of the blown-up
$X_0$ is extracted by Hurwitz moves as a prefix of model words from
$E(g+1)$, which have nonzero extremal coefficient. Then apply the KM structure
theorem and the blow-up formula.
Insert $T^k$ at $Y_0$. Surface insertions crossing $T$ are tracked by a
nilpotent deformation over $\Q[t_i]/(t_i^2)$, and the deformed $T_L$ is still
a quasi-isomorphism. Cayley--Hamilton gives $k\ge1$ with nonzero all-label
coefficient. Cut at an internal $Y_2$ seam, where $H_2$ splits over $\Q$, and
allocate the probes. This gives caps $C_\pm$ with $b^+>0$, containing $A$
with $w\cdot A$ odd, so generic periods kill reducibles. Each cap gets one
extra blow-up $E_{\rm exc}$. The result is an integral cycle $x$ and cocycle
$y$ on $I(Y_2)$ with $yx=q\ne0$.

\section*{3. The stack: $\Omega\ne0$ and $n_D>0$ (paper, Sec.~5)}
$X=C_-\cup(\text{$n$ copies of }B,\ \text{each followed by }\phi)\cup C_+$,
except at $m$ places where $H=f_+f_-$ replaces the return. These positive
bridges are spaced exactly four intervals apart.
\begin{itemize}
\item \emph{Integrality.} Each step is a unit of $L/2^NL$, where
$L=I(Y_2;\Z)/\mathrm{Tor}$ and $q\not\equiv0\ (2^N)$. Take all repetition
counts divisible by $\exp\mathrm{GL}_r(\Z/2^N)$ and by $8$. Then the composite
is the identity mod $2^N$, so $\Omega\equiv\pm q\ne0$. Stretching the finite
$Y$-seams identifies $\Omega$ with the composite: a central contact costs
$3$, and a macro $BHB$ recovers at most $2$ per central outer limit.
\item \emph{Lattices.} Cut at the $S^3$'s. A negative cell is unimodular,
$\langle a,b\rangle$ with $a^2=b^2=-1$. A positive cell has
$H_2=\bigl(\begin{smallmatrix}2&1\\1&0\end{smallmatrix}\bigr)\oplus(-1)^4$,
$b^+=1$, $\sigma=-4$, and contains the plumbing
$X_1=-S_j,\ U,\ X_2=S_{j+1}$ with squares $(-4,0,-4)$.
\item \emph{Indices.} Fix a characteristic $l$ and put $\Lambda=l+c$. A bit
flip adds $\PD(S_i)$ to both $c$ and $\Lambda$, preserving $\Lambda^2$ and
$\kappa$. $\Theta=(\Lambda^2-\sigma)/4$ telescopes, gaining $+1$ per positive
cell. With $D_I=8\kappa-3(1+b^+)=n+z$ and $b^+=m+b_0$,
$n_D=\Theta-\kappa=\frac{5m-n}{8}+\mathrm{const}$, which is about $m/8$ for $n\approx4m$.
\end{itemize}
\emph{Where $\phi$ enters.} $n_D>0$ needs $m/n>1/5$. The exclusions of
part~5 below need spacing $\ge4$, i.e.\ $m/n\le1/4$. Without $\phi$ every $B$
must be followed by $H$, since the reversed interval
$Y_{-2}\to Y_2$ has $b^+=2$. That forces $m\approx n$.

\section*{4. Metrics and the Abelian clamp (paper, Secs.~6--7)}
Only the $S^3$ cuts $J_i$ and the lens cuts $\partial N_i$ become infinite;
both have psc. All $Y$-seams stay finite, so no monopole Floer theory on
$Y_j$ is needed. In each macro the two adjacent interval parameters drive
a toric K\"ahler metric (Guillemin) on the $(-4,0,-4)$ plumbing. Its moment
polygon is $\{4|y|\le x\le1,\,-a\le y\le b\}$, its class is
$H=U+aX_1+bX_2$, and $0\le a,b\le\frac14$. The value $a=\frac14$ splits off
$N_1$ (the lens face); $a\to0$ approaches $J$. Abreu's formula plus a Schur
product identity give $\mathrm{Scal}\ge0$. Near $a=b=0$ a round
$S^2\times S^1$ tube in class $U$ has length $\sim1/\epsilon$.
For reductions (R: $\Phi=0$ split; S: $\Phi$ in one summand) let $v$ be the
line difference, $K=\Lambda+v$ and $u=v\cdot U$.
\begin{itemize}
\item Small width: the energy is $o(1/\epsilon)$, but $u\ne0$ costs
$\gtrsim u^2/\epsilon$. Hence $u=0$.
\item Positive width: Weitzenb\"ock gives $vH=0$ (R) or $(vH)(KH)<0$ (S).
$H$ is a convex combination of the vertex classes
$H_I=U+\frac14\sum_{i\in I}X_i$, and $\Lambda H_I<0$. So for some
$I\ne\emptyset$: $\operatorname{sign}(u)\,vH_I\le0$ (R) or
$<|\Lambda H_I|$ (S).
\item Caps: no R component meets a cap. S components satisfy
$v_W^2,\ K_W^2\le C_W$, with $C_W$ independent of the stack.
\end{itemize}

\section*{5. Index exclusion (paper, Sec.~8)}
The equations are $D_a\Phi=0$, $\rho(F_a^{0,+})=(\Phi\Phi^*)_{00}$, with
determinant-one gauge and $S^1$ acting on $\Phi$. Fixed types are A
(instanton), R and S; free type is P. Take a limit with $k$ cuts at the
$J_i$. Virtually fill the lens caps with sharp references; the framed lens
corrections come from $\mathbb F_4$ versus $\mathbb P(1,1,4)$. Then
$q_i=D_{I,i}+p_i-2s_i-z_i$ satisfies $\sum(q_i+4)=4$.
On R/S components $\kappa_i=-v_i^2/4+\ell_i$ with $\ell_i\ge0$. A test with
$v(S)=0$ on a reducible forces a unit of $\ell$; this uses a smooth
zero-winding modification of the holonomy test. Every assigned test pays a
negative half of square $\le-2$ or a unit of $\ell$. Every positive cell has
square $\le-2$ by the clamp, with one exception that is paid by $\ell$.
With $L\ge4m-3$ (spacing four), every interior run of R/S components has
$\sum(q_i+4)\ge-2$. A cap component is never R. An S cap component has
$q+4\ge8$ after long delays and large odd $\Lambda(E_{\rm exc})$, using both
cap bounds. So $k\ge1$ would force $\sum(q_i+4)\ge2a+6b+2\ge6>4$, and
$k=0$ forces type A. \emph{Hence there are no Seiberg--Witten (S) or R fixed
points at all.} For a mixed limit, delete the R components and keep their
parameters. For $k\ge1$ there are $N\ge2$ non-R components and the projected
dimension is $\le3-2N<0$. The only survivor is one P
component plus one split $N_i$ with the minimal trace cap ($\kappa=\frac14$)
and no particles.

\section*{6. Boundary count (paper, Secs.~9--10)}
Instanton links: make the Dirac operator surjective; the link is
$\CP^{n_D-1}$ and the phase cuts are $\mathcal O(2)$-sections, giving
$2^\eta\operatorname{sign}$. Lens ends: on $N$ the spinor vanishes and the
Dirac operator is invertible (psc, index $0$), and the cap glues to the free
main solution. The two bit lifts contribute opposite terms with the same
$\varepsilon$ as in $B$; the Dirac line is complex, so it adds no sign.
Stokes (rational multisections on the free stratum) gives
$0=\pm2^\eta\Omega$, a contradiction.\hfill$\square$

\section*{Assessment}
The $\mathrm{PU}(2)$-monopole cobordism was built to relate Donaldson and
Seiberg--Witten invariants. Using it as a vanishing mechanism for a
$3$-manifold problem is new. Here the geometry empties the Seiberg--Witten
side, so the Feehan--Leness reducible gluing is never needed. Also new:
(i)~the integral negative-interval map $B$, whose lens-wall cancellation
carries over to the monopole problem; (ii)~using $\phi$ to set the
positive-bridge density in $(\frac15,\frac14]$; (iii)~the toric psc families
and the period clamp on the $b^+=1$ cells. The mod-$2^N$ repetition and the
labelled deformation are routine. The rest of the length is index
bookkeeping and analytic verification. Points to check, in my reading:
\begin{itemize}
\item transversality and compactness of the ``necessary projection'' that
forgets R components in mixed strata, with multisections;
\item the zero-winding holonomy test: it must preserve the class of $B$ and
behave correctly under Uhlenbeck limits;
\item uniformity of the clamp across the degenerating toric family and
genuine faces, and its compatibility with the product-metric computation of
$\Omega$;
\item the local orientation comparison that cancels lens ends in both
problems.
\end{itemize}
\end{document}
